% Plantilla — Dossier PDF de problemas de investigación (skill: proponer-problemas-pdf)
% Compila con: pdflatex -interaction=nonstopmode -halt-on-error doc.tex  (2 pasadas)
% Sustituye los <<MARCADORES>> y los placeholders {{...}}. No añadas paquetes exóticos: este preámbulo compila en TeX Live estándar.
\documentclass[11pt]{article}
\usepackage[utf8]{inputenc}
\usepackage[T1]{fontenc}
\usepackage[spanish,es-noindentfirst]{babel}
\usepackage[a4paper,margin=2.4cm]{geometry}
\usepackage{amsmath,amssymb,amsthm}
\usepackage{enumitem}
\usepackage{xcolor}
\usepackage{mdframed}
\usepackage{booktabs}
\usepackage[hidelinks]{hyperref}
\usepackage{microtype}

\definecolor{acc}{RGB}{30,80,140}
\definecolor{soft}{RGB}{245,247,250}

\newtheorem{problema}{Problema}
\newtheorem*{nota}{Nota}

% Caja resaltada (para "hueco unificador", advertencias, etc.)
\newmdenv[backgroundcolor=soft,linecolor=acc,linewidth=0.6pt,
  roundcorner=3pt,skipabove=6pt,skipbelow=6pt,
  innertopmargin=6pt,innerbottommargin=6pt]{caja}

% Atajos frecuentes (ajústalos al dominio si hace falta)
\newcommand{\R}{\mathbb{R}}
\newcommand{\Z}{\mathbb{Z}}
\newcommand{\N}{\mathbb{N}}

\title{\vspace{-1.2cm}\textbf{<<TÍTULO>>}\\[4pt]
\large <<SUBTÍTULO: p. ej. Fuentes y problemas de investigación en {{DOMINIO}}>>}
\author{Compilación para el programa de {{DOMINIO}}}
\date{<<FECHA>>}

\begin{document}
\maketitle
\vspace{-0.6cm}

\begin{abstract}
\noindent <<ABSTRACT: de qué corpus se parte, qué se sintetiza, énfasis (líneas del dominio), y RB de respaldo si lo hay.>>
\end{abstract}

\section{Marco mínimo}
<<Definiciones/contexto imprescindibles del tema, concisos, con su reducción al caso base.>>

\section{Fuentes interesantes}
\subsection*{Fundacionales}
\begin{itemize}[leftmargin=1.4em,topsep=2pt,itemsep=3pt]
  \item \textbf{<<Autor (Fecha)>>.} <<aporte central + problema abierto/limitación (cita textual si la hay).>>
\end{itemize}
\subsection*{Modernas / aplicaciones}
\begin{itemize}[leftmargin=1.4em,topsep=2pt,itemsep=3pt]
  \item \textbf{<<Autor (Fecha)>>.} <<resultado clave; qué cubre y qué NO.>>
\end{itemize}

\begin{caja}
\textbf{Lectura conjunta (el hueco unificador).} <<Cruza qué cubre cada fuente y cuál es la intersección que ninguna cubre — de ahí salen los problemas.>>
\end{caja}

\section{Problemas de investigación propuestos}
\subsection*{<<Grupo A: p. ej. Fundamentos>>}
\begin{problema}[<<título corto>>]
<<Enunciado preciso. \emph{Base:} <<fuente>>. \emph{Hueco:} <<qué falta>>.>>
\end{problema}

\subsection*{<<Grupo B>>}
\begin{problema}[<<título corto>>]
<<...>>
\end{problema}

\section{Priorización y advertencias}
\begin{center}\small
\begin{tabular}{@{}p{5.6cm}cccc@{}}
\toprule
\textbf{Problema} & \textbf{Tipo} & \textbf{Base} & \textbf{Novedad} & \textbf{Viabilidad}\\
\midrule
<<P1>> & <<...>> & <<...>> & <<...>> & <<...>>\\
\bottomrule
\end{tabular}
\end{center}

\begin{caja}
\textbf{Advertencias de prioridad.} <<Qué podría colisionar; recomendar cierre de prioridad con el agente \emph{investigador} antes de redactar. Distinguir problemas abiertos de autor vs. propuestas propias.>>
\end{caja}

\begin{thebibliography}{99}\small
\bibitem{ref1} <<Apellido, N. (Fecha). \emph{Título}. Venue vol(núm), págs.>>
\end{thebibliography}

\end{document}
